% GENERATED FILE. Edit canonical lessons, then run npm run generate:books.
% canonical-source-hash: fnv1a64:e85375e45f087e97
% canonical-lessons: GU-C256-citra, GU-C256-citarvu, GU-C256-pichi, GU-C256-photo, GU-C256-guthvu

\chapter{Paint, photos and knitting}
\label{ch:gu-paint-photos-and-knitting}
\hlchaptermodality{256}

\begin{chapteropening}
\textbf{By the end of this chapter:} \emph{I can say what I painted, photographed or knitted.}

Complete the last lesson of chapter 256: I can say what I painted, photographed or knitted.
\end{chapteropening}

\section[\texorpdfstring{\gu{ચિત્ર} (citra)}{ચિત્ર (citra)}]{\gu{ચિત્ર} (citra) — picture, painting}

\label{lesson:GU-C256-citra}



\begin{quote}
Before the new one: say the Gujarati for drama, then the Gujarati for film.
\end{quote}

\subsection*{You'll want to know: \gu{ચિત્ર}}
\textbf{\gu{ચિત્ર}} — \emph{citra} — \textquotedblleft{}picture, painting\textquotedblright{}.

\begin{quote}
\textbf{\gu{બાળકે} \gu{સુંદર} \gu{ચિત્ર} \gu{બનાવ્યું}.} — \emph{bāḷake sundar citra banāvyũ} — \textquotedblleft{}The child made a beautiful picture.\textquotedblright{}
\end{quote}

\begin{cousinweb}
\textbf{\gu{ચિત્ર}} (\emph{citra}) is Sanskrit \emph{citrá}, \textquotedblleft{}conspicuous, bright, variegated\textquotedblright{}, taken over as a tatsama for a picture (CDIAL 4803).
\end{cousinweb}

\begin{grammarlens}[title={What you've built}]
One more word for everyday life at A2.
\end{grammarlens}

\subsection*{Your turn}
\begin{itemize}
\raggedright
  \item \emph{Say it:} \emph{citra}
  \item \emph{Say it:} \emph{citra}, once more
  \item \emph{Say it:} \emph{bāḷake sundar citra banāvyũ} — The child made a beautiful picture.
  \item \emph{Recall:} say the Gujarati for drama, then the Gujarati for film, then say \emph{citra} again
\end{itemize}

\begin{tcolorbox}[breakable,colback=teal!4,colframe=teal!35!black,title={Before you move on}]
What does \gu{ચિત્ર} mean? (\textquotedblleft{}Picture, painting\textquotedblright{}.) Say it once more.
\end{tcolorbox}
\section[\texorpdfstring{\gu{ચીતરવું} (cītarvũ)}{ચીતરવું (cītarvũ)}]{\gu{ચીતરવું} (cītarvũ) — to paint, to make pictures}

\label{lesson:GU-C256-citarvu}



\begin{quote}
Before the new one: say the Gujarati for film, then the Gujarati for picture.
\end{quote}

\subsection*{You'll want to know: \gu{ચીતરવું}}
\textbf{\gu{ચીતરવું}} — \emph{cītarvũ} — \textquotedblleft{}to paint, to make pictures\textquotedblright{}.

\begin{quote}
\textbf{\gu{તેણે} \gu{એક} \gu{મોર} \gu{ચીતર્યો}.} — \emph{teṇe ek mor cītaryo} — \textquotedblleft{}She painted a peacock.\textquotedblright{}
\end{quote}

\begin{cousinweb}
\textbf{\gu{ચીતરવું}} (\emph{cītarvũ}) is a tadbhava. It comes from Sanskrit \emph{citráyati}, \textquotedblleft{}decorates\textquotedblright{}, built on \emph{citrá}, \textquotedblleft{}bright\textquotedblright{}; CDIAL lists Gujarati \emph{citarvũ} (CDIAL 4810).
\end{cousinweb}

\begin{grammarlens}[title={What you've built}]
One more word for everyday life at A2.
\end{grammarlens}

\subsection*{Your turn}
\begin{itemize}
\raggedright
  \item \emph{Say it:} \emph{cītarvũ}
  \item \emph{Say it:} \emph{cītarvũ}, once more
  \item \emph{Say it:} \emph{teṇe ek mor cītaryo} — She painted a peacock.
  \item \emph{Recall:} say the Gujarati for film, then the Gujarati for picture, then say \emph{cītarvũ} again
\end{itemize}

\begin{tcolorbox}[breakable,colback=teal!4,colframe=teal!35!black,title={Before you move on}]
What does \gu{ચીતરવું} mean? (\textquotedblleft{}To paint, to make pictures\textquotedblright{}.) Say it once more.
\end{tcolorbox}
\section[\texorpdfstring{\gu{પીંછી} (pī̃chī)}{પીંછી (pī̃chī)}]{\gu{પીંછી} (pī̃chī) — paintbrush}

\label{lesson:GU-C256-pichi}



\begin{quote}
Before the new one: say the Gujarati for picture, then the Gujarati for to paint.
\end{quote}

\subsection*{You'll want to know: \gu{પીંછી}}
\textbf{\gu{પીંછી}} — \emph{pī̃chī} — \textquotedblleft{}paintbrush\textquotedblright{}.

\begin{quote}
\textbf{\gu{મેં} \gu{નવી} \gu{પીંછી} \gu{ખરીદી}.} — \emph{mẽ navī pī̃chī kharīdī} — \textquotedblleft{}I bought a new paintbrush.\textquotedblright{}
\end{quote}

\begin{cousinweb}
\textbf{\gu{પીંછી}} (\emph{pī̃chī}), a little feather, is built on \emph{pī̃ch}, \textquotedblleft{}feather\textquotedblright{}, with \emph{-ī}. CDIAL derives \emph{pī̃ch} from Sanskrit \emph{piccha}, \textquotedblleft{}tail-feather\textquotedblright{} (CDIAL 8151), but lists no \emph{pī̃chī}.
\end{cousinweb}

\begin{grammarlens}[title={What you've built}]
One more word for everyday life at A2.
\end{grammarlens}

\subsection*{Your turn}
\begin{itemize}
\raggedright
  \item \emph{Say it:} \emph{pī̃chī}
  \item \emph{Say it:} \emph{pī̃chī}, once more
  \item \emph{Say it:} \emph{mẽ navī pī̃chī kharīdī} — I bought a new paintbrush.
  \item \emph{Recall:} say the Gujarati for picture, then the Gujarati for to paint, then say \emph{pī̃chī} again
\end{itemize}

\begin{tcolorbox}[breakable,colback=teal!4,colframe=teal!35!black,title={Before you move on}]
What does \gu{પીંછી} mean? (\textquotedblleft{}Paintbrush\textquotedblright{}.) Say it once more.
\end{tcolorbox}
\section[\texorpdfstring{\gu{ફોટો} (phoṭo)}{ફોટો (phoṭo)}]{\gu{ફોટો} (phoṭo) — photo}

\label{lesson:GU-C256-photo}



\begin{quote}
Before the new one: say the Gujarati for to paint, then the Gujarati for paintbrush.
\end{quote}

\subsection*{You'll want to know: \gu{ફોટો}}
\textbf{\gu{ફોટો}} — \emph{phoṭo} — \textquotedblleft{}photo\textquotedblright{}.

\begin{quote}
\textbf{\gu{મેં} \gu{દરિયાનો} \gu{ફોટો} \gu{લીધો}.} — \emph{mẽ dariyāno phoṭo līdho} — \textquotedblleft{}I took a photo of the sea.\textquotedblright{}
\end{quote}

\begin{cousinweb}
\textbf{\gu{ફોટો}} (\emph{phoṭo}) is a loan from English \emph{photo}.
\end{cousinweb}

\begin{grammarlens}[title={What you've built}]
One more word for everyday life at A2.
\end{grammarlens}

\subsection*{Your turn}
\begin{itemize}
\raggedright
  \item \emph{Say it:} \emph{phoṭo}
  \item \emph{Say it:} \emph{phoṭo}, once more
  \item \emph{Say it:} \emph{mẽ dariyāno phoṭo līdho} — I took a photo of the sea.
  \item \emph{Recall:} say the Gujarati for to paint, then the Gujarati for paintbrush, then say \emph{phoṭo} again
\end{itemize}

\begin{tcolorbox}[breakable,colback=teal!4,colframe=teal!35!black,title={Before you move on}]
What does \gu{ફોટો} mean? (\textquotedblleft{}Photo\textquotedblright{}.) Say it once more.
\end{tcolorbox}
\section[\texorpdfstring{\gu{ગૂંથવું} (gū̃thvũ)}{ગૂંથવું (gū̃thvũ)}]{\gu{ગૂંથવું} (gū̃thvũ) — to knit, to plait}

\label{lesson:GU-C256-guthvu}



\begin{quote}
Before the new one: say the Gujarati for paintbrush, then the Gujarati for photo.
\end{quote}

\subsection*{You'll want to know: \gu{ગૂંથવું}}
\textbf{\gu{ગૂંથવું}} — \emph{gū̃thvũ} — \textquotedblleft{}to knit, to plait\textquotedblright{}.

\begin{quote}
\textbf{\gu{મારી} \gu{દાદીએ} \gu{સ્વેટર} \gu{ગૂંથ્યું}.} — \emph{mārī dādīe sveṭar gū̃thyũ} — \textquotedblleft{}My grandmother knitted a sweater.\textquotedblright{}
\end{quote}

\begin{cousinweb}
\textbf{\gu{ગૂંથવું}} (\emph{gū̃thvũ}) is a tadbhava. CDIAL lists it, glossed \textquotedblleft{}to knit, plait\textquotedblright{}, under Sanskrit \emph{gupháti}, \textquotedblleft{}strings together\textquotedblright{}, from a doubtful participle \emph{*guptha-} blended with the Prakrit \emph{gaṁthaï}, \textquotedblleft{}ties\textquotedblright{} (CDIAL 4205).
\end{cousinweb}

\begin{grammarlens}[title={What you've built}]
One more word for everyday life at A2.
\end{grammarlens}

\subsection*{Your turn}
\begin{itemize}
\raggedright
  \item \emph{Say it:} \emph{gū̃thvũ}
  \item \emph{Say it:} \emph{gū̃thvũ}, once more
  \item \emph{Say it:} \emph{mārī dādīe sveṭar gū̃thyũ} — My grandmother knitted a sweater.
  \item \emph{Recall:} say the Gujarati for paintbrush, then the Gujarati for photo, then say \emph{gū̃thvũ} again
\end{itemize}

\begin{tcolorbox}[breakable,colback=teal!4,colframe=teal!35!black,title={Before you move on}]
What does \gu{ગૂંથવું} mean? (\textquotedblleft{}To knit, to plait\textquotedblright{}.) Say it once more.
\end{tcolorbox}
